% Unit 7 (Equilibrium) guided notes -- 8 blocks, CED sequence.
\documentclass{apchem-notes}

\apcunit{7}
\apcunittitle{Equilibrium}
\apcdoctitle{Guided Notes}
\apcsubtitle{CED 7.1--7.12 \textbullet\ Zumdahl Ch 13, \S16.1 \textbullet\ 8 blocks}

\begin{document}
\apctitleblock

\begin{apcnote}[The comparison this whole unit runs on]
One question answers almost everything in this unit:
\textbf{how does $Q$ compare with $K$?}

\begin{center}
\begin{tabular}{@{}lll@{}}
\hline
$Q < K$ & too few products & shifts \textbf{right} \\
$Q = K$ & at equilibrium & no net change \\
$Q > K$ & too many products & shifts \textbf{left} \\
\hline
\end{tabular}
\end{center}

That single comparison tells you which way a fresh mixture will move,
whether a system has arrived, what a disturbance will do, and --- in the
last block --- whether a precipitate forms. Le Ch\^atelier's principle is a
shortcut for it, not a separate idea.
\end{apcnote}

% =====================================================================
\blocklesson{Dynamic Equilibrium}{CED 7.1--7.2 \textbullet\ Zumdahl \S13.1}

\begin{objectives}
  \item Describe equilibrium as equal forward and reverse rates.
  \item Interpret concentration-time graphs approaching equilibrium.
\end{objectives}

\reading{Zumdahl \S13.1}{651--654}

\begin{warmup}[4]
  \item A reversible reaction is written with what symbol?
        \answerblank[2.6cm]{$\rightleftharpoons$}
  \item Rate depends on: \answerblank[3cm]{concentration}
  \item As reactant is consumed, the forward rate:
        \answerblank[2.4cm]{decreases}
  \item As product builds up, the reverse rate:
        \answerblank[2.4cm]{increases}
\end{warmup}

\segment{INSTRUCTION A}{25}{What ``equilibrium'' actually means}

\notesection{Equal rates, not equal amounts}{7.1}
\practice{6}

Start with pure reactants. The forward rate is at its
\answerblank[2cm]{maximum} and the reverse rate is
\answerblank[1.6cm]{zero}. As the reaction proceeds:

\begin{itemize}[leftmargin=*,itemsep=0.3em]
  \item Reactant concentration falls, so the forward rate
        \answerblank[2.2cm]{falls}.
  \item Product concentration rises, so the reverse rate
        \answerblank[2.2cm]{rises}.
  \item When the two rates become \answerblank[2cm]{equal}, concentrations
        stop changing. That is \term{dynamic equilibrium}.
\end{itemize}

\begin{apctrap}
\textbf{Equilibrium does not mean equal concentrations.} It means equal
\emph{rates}. The concentrations at equilibrium are almost never equal ---
their ratio is fixed by $K$, which can be $10^{-10}$ or $10^{10}$.

It also does not mean the reaction has stopped. Both directions continue at
the same rate, which is why the word is \emph{dynamic}. A question asking
``what is happening at equilibrium?'' wants that sentence.
\end{apctrap}

\begin{center}
\begin{tikzpicture}
\begin{axis}[width=10.2cm, height=4.8cm,
    xlabel={time}, ylabel={concentration},
    xtick=\empty, ytick=\empty, xmin=0, xmax=10, ymin=0, ymax=10,
    axis lines=left, legend style={draw=none, font=\sffamily\scriptsize,
      at={(0.98,0.6)}, anchor=east}]
  \addplot[seriesA, samples=120, domain=0:10] {2 + 6*exp(-0.55*x)};
  \addlegendentry{reactant}
  \addplot[seriesB, samples=120, domain=0:10] {8 - 6*exp(-0.55*x)};
  \addlegendentry{product}
  \draw[densely dotted] (axis cs:6,0) -- (axis cs:6,8.5);
  \node[font=\scriptsize,anchor=south] at (axis cs:7.6,8.6)
    {equilibrium reached};
\end{axis}
\end{tikzpicture}
\end{center}

Both curves \answerblank[2.4cm]{level off} --- and they level off at
\emph{different} heights, which is the graphical version of the trap above.

\segment{APPLICATION}{20}{Reading the graph}

\begin{enumerate}[leftmargin=*,itemsep=0.4em]
  \item At the moment equilibrium is reached, compare the forward and
        reverse rates, and compare the concentrations.
        \answerlines[3]{The rates are equal --- that is the definition. The
        concentrations are generally unequal and simply stop changing;
        their ratio is whatever $K$ requires.}
  \item A student says ``the reaction stopped.'' Correct them.
        \answerlines[2]{Nothing stopped. Both reactions continue at equal
        rates, so every product molecule formed is matched by one
        converting back, and no net change is observable.}
  \item How would the graph differ for a reaction with a very large $K$?
        \answerlines[2]{The product curve would level off much higher and
        the reactant curve much lower --- the plateau heights encode the
        magnitude of $K$.}
\end{enumerate}

\begin{exitticket}
Give two things that are equal at equilibrium and one thing that generally
is not.
\keyonly{\begin{apcanswer}Equal: the forward and reverse rates, and (by
consequence) the rates of change of every concentration, which are all
zero. Not equal: the concentrations themselves, whose ratio is set by
$K$.\end{apcanswer}}
\end{exitticket}

% =====================================================================
\blocklesson{$Q$, $K$, and Calculating Them}{CED 7.3--7.4 \textbullet\ Zumdahl \S13.2--13.4}

\begin{objectives}
  \item Write $K$ and $Q$ expressions correctly.
  \item Calculate $K$ from equilibrium data.
\end{objectives}

\reading{Zumdahl \S13.2--13.4}{654--662}

\begin{warmup}[3]
  \item Equilibrium means equal: \answerblank[2cm]{rates}
  \item Coefficients become what in a $K$ expression?
        \answerblank[2.6cm]{exponents}
  \item Which species are omitted from $K$?
        \answerblank[4.8cm]{pure solids and liquids}
\end{warmup}

\segment{INSTRUCTION A}{25}{Writing the expression}

\notesection{$K$ and $Q$ have the same form}{7.3}
\practice{5}

For $\ce{aA + bB <=> cC + dD}$:
\[ K = \frac{[\ce{C}]^c[\ce{D}]^d}{[\ce{A}]^a[\ce{B}]^b} \]

$Q$ has the \emph{identical} form but uses concentrations at
\answerblank[3.4cm]{any moment}, not necessarily equilibrium ones.

\begin{apctrap}
\textbf{Omit pure solids and pure liquids.} Their ``concentration'' does not
vary --- a bigger lump of solid is not more concentrated --- so they carry
no term. Aqueous species and gases always appear.

Watch for water: \ce{H2O(l)} as a solvent is omitted, but \ce{H2O(g)} in a
gas-phase reaction is \emph{included}. The state symbol is doing real work
and is not decoration.
\end{apctrap}

\begin{workedexample}[Worked example 1: expressions]
\begin{align*}
  \ce{N2(g) + 3H2(g) &<=> 2NH3(g)}
    &K &= \frac{[\ce{NH3}]^2}{[\ce{N2}][\ce{H2}]^3} \\[0.3em]
  \ce{CaCO3(s) &<=> CaO(s) + CO2(g)}
    &K &= [\ce{CO2}] \\[0.3em]
  \ce{2H2O(l) &<=> 2H2(g) + O2(g)}
    &K &= [\ce{H2}]^2[\ce{O2}]
\end{align*}
The second has only one term --- both solids drop out. The third omits
liquid water but would keep it if it were steam.
\end{workedexample}

\segment{INSTRUCTION B}{20}{Calculating $K$}

\notesection{From equilibrium concentrations}{7.4}
\practice{5}

\begin{workedexample}[Worked example 2: substituting into $K$]
At equilibrium a vessel contains $[\ce{N2}] = 0.50$,
$[\ce{H2}] = 0.30$, $[\ce{NH3}] = 0.20$~M for
$\ce{N2 + 3H2 <=> 2NH3}$.

\[ K = \frac{(0.20)^2}{(0.50)(0.30)^3}
     = \frac{0.040}{(0.50)(0.027)} = \frac{0.040}{0.0135}
     = \mathbf{3.0} \]

The cube on hydrogen is where marks are lost --- $(0.30)^3 = 0.027$, not
$0.90$.
\end{workedexample}

\notesection{$K_p$ and $K_c$}{7.4}

For gases, $K$ may be written with partial pressures:
\[ K_p = K_c(RT)^{\Delta n} \qquad
   \Delta n = \text{mol gas products} - \text{mol gas reactants} \]

When $\Delta n = 0$, $K_p$ and $K_c$ are \answerblank[2cm]{equal}.

\begin{workedexample}[Worked example 3: converting]
For $\ce{N2 + 3H2 <=> 2NH3}$ at \SI{500}{\kelvin} with
$K_c = 0.50$: $\Delta n = 2 - 4 = -2$.
\[ K_p = (0.50)[(0.08206)(500)]^{-2} = \frac{0.50}{(41.03)^2}
   = \mathbf{3.0\times10^{-4}} \]
\end{workedexample}

\segment{APPLICATION}{20}{Expressions and values}

\begin{enumerate}[leftmargin=*,itemsep=0.4em]
  \item Write $K$ for $\ce{2SO2(g) + O2(g) <=> 2SO3(g)}$.
        \answerblank[6cm]{$[\ce{SO3}]^2/([\ce{SO2}]^2[\ce{O2}])$}
  \item Write $K$ for $\ce{Fe3O4(s) + 4H2(g) <=> 3Fe(s) + 4H2O(g)}$.
        \answerblank[6cm]{$[\ce{H2O}]^4/[\ce{H2}]^4$}
  \item At equilibrium $[\ce{H2}] = 0.20$, $[\ce{I2}] = 0.40$,
        $[\ce{HI}] = 0.60$~M. Find $K$ for
        $\ce{H2 + I2 <=> 2HI}$. \workspace[2cm]
        \keyonly{\begin{apcanswer}$K = (0.60)^2/[(0.20)(0.40)] =
        0.36/0.080 = \mathbf{4.5}$\end{apcanswer}}
\end{enumerate}

\begin{exitticket}
Why do pure solids and liquids not appear in $K$, and what determines
whether water appears?
\keyonly{\begin{apcanswer}Their concentration is a fixed property of the
substance and does not change as the reaction proceeds, so including them
would add a constant. Water appears when it is a gas or a dissolved
reactant, and is omitted when it is the pure liquid solvent --- the state
symbol decides.\end{apcanswer}}
\end{exitticket}

% =====================================================================
\blocklesson{Magnitude and Properties of $K$}{CED 7.5--7.6 \textbullet\ Zumdahl \S13.3--13.5}

\begin{objectives}
  \item Interpret the size of $K$.
  \item Adjust $K$ when a reaction is manipulated.
\end{objectives}

\reading{Zumdahl \S13.3--13.5}{658--671}

\begin{warmup}[3]
  \item $K$ for $\ce{CaCO3(s) <=> CaO(s) + CO2(g)}$:
        \answerblank[2.6cm]{$[\ce{CO2}]$}
  \item $K_p = K_c(RT)^{?}$: \answerblank[2cm]{$\Delta n$}
  \item Only what change alters the value of $K$?
        \answerblank[2.6cm]{temperature}
\end{warmup}

\segment{INSTRUCTION A}{25}{What the number tells you}

\notesection{Reading the magnitude}{7.5}
\practice{5}

\begin{center}
\begin{tabular}{@{}ll@{}}
\hline
$K \gg 1$ (say $> 10^2$) & \textbf{products} strongly favored \\
$K \approx 1$ & appreciable amounts of both \\
$K \ll 1$ (say $< 10^{-2}$) & \textbf{reactants} strongly favored \\
\hline
\end{tabular}
\end{center}

\begin{apctrap}
$K$ says \emph{how far}, never \emph{how fast}. A reaction with
$K = 10^{40}$ can be immeasurably slow --- that is the kinetic control you
met in Unit 5. The two questions are independent, and the exam likes pairing
them to see whether you have separated them.
\end{apctrap}

\segment{INSTRUCTION B}{20}{Manipulating $K$}

\notesection{Four rules}{7.6}
\practice{5}

\begin{center}
\begin{tabular}{@{}ll@{}}
\hline
\textbf{Change to the equation} & \textbf{New $K$} \\
\hline
reverse it & $1/K$ \\
multiply the coefficients by $n$ & $K^n$ \\
divide the coefficients by $n$ & $K^{1/n}$ \\
add two reactions & $K_1 \times K_2$ \\
\hline
\end{tabular}
\end{center}

\begin{workedexample}[Worked example 4: one reaction, four values]
If $K = 4.0$ for $\ce{A + B <=> C}$, then:

\begin{itemize}[leftmargin=*,itemsep=0.2em]
  \item $\ce{C <=> A + B}$ has $K = 1/4.0 = \mathbf{0.25}$.
  \item $\ce{2A + 2B <=> 2C}$ has $K = (4.0)^2 = \mathbf{16}$.
  \item $\ce{\tfrac{1}{2}A + \tfrac{1}{2}B <=> \tfrac{1}{2}C}$ has
        $K = \sqrt{4.0} = \mathbf{2.0}$.
  \item Added to a second reaction with $K = 3.0$, the overall $K$ is
        $(4.0)(3.0) = \mathbf{12}$.
\end{itemize}

Notice these are \emph{multiplicative} operations, unlike $\Delta H$ in
Unit 6, which was additive. Equilibrium constants multiply where enthalpies
add.
\end{workedexample}

\segment{APPLICATION}{20}{Using the rules}

\begin{enumerate}[leftmargin=*,itemsep=0.4em]
  \item $K = 2.5\times10^{-3}$ for a reaction. Give $K$ for the reverse.
        \answerblank[3cm]{$4.0\times10^{2}$}
  \item Classify each as products-favored, reactants-favored, or mixed:
        $K = 1\times10^{-10}$, $K = 1.5$, $K = 3\times10^{6}$.
        \answerlines[2]{Reactants strongly favored; appreciable amounts of
        both; products strongly favored.}
  \item A reaction has $K = 10^{35}$ but no observable change occurs when
        the reactants are mixed. Is the value wrong? Explain.
        \answerlines[3]{No. $K$ describes the position of equilibrium, not
        the rate of reaching it. A very large activation energy can make the
        reaction immeasurably slow --- it is under kinetic control, and the
        thermodynamic prediction is untouched.}
\end{enumerate}

\begin{exitticket}
State what happens to $K$ when a reaction is reversed and when its
coefficients are doubled, and name the one change that alters $K$ for a
given reaction.
\keyonly{\begin{apcanswer}Reversing gives $1/K$; doubling the coefficients
gives $K^2$. Only a change in \emph{temperature} alters $K$ for a given
reaction as written --- concentration, pressure and catalysts do
not.\end{apcanswer}}
\end{exitticket}

% =====================================================================
\blocklesson{ICE Tables}{CED 7.7 \textbullet\ Zumdahl \S13.6}

\begin{objectives}
  \item Set up and solve an ICE table.
  \item Apply and test the small-$x$ approximation.
\end{objectives}

\reading{Zumdahl \S13.6}{671--676}

\begin{warmup}[3]
  \item Reversing a reaction gives $K =$ \answerblank[1.8cm]{$1/K$}
  \item $K \ll 1$ favors: \answerblank[2.4cm]{reactants}
  \item BCA tables from Unit 4 tracked what?
        \answerblank[4.2cm]{moles, to completion}
\end{warmup}

\segment{INSTRUCTION A}{25}{The table you already know}

\notesection{ICE is BCA with an unknown}{7.7}
\practice{5}

\term{ICE} stands for \answerblank[5.8cm]{Initial, Change, Equilibrium}. It
is the Before/Change/After table from stoichiometry, with one difference:
the reaction stops partway, so the change is an unknown
\answerblank[1.4cm]{$x$}.

The change row is still driven by the \answerblank[2.6cm]{coefficients} ---
if A has coefficient 1 and C has coefficient 2, then A changes by $-x$ and
C by $+2x$.

\begin{workedexample}[Worked example 5: the small-$x$ approximation]
$\ce{A <=> B + C}$ with $K = 1.0\times10^{-5}$ and
$[\ce{A}]_0 = \SI{0.100}{\Molar}$.

\begin{center}
\begin{tabular}{@{}lccc@{}}
\hline
 & \ce{A} & \ce{B} & \ce{C} \\
\hline
Initial & 0.100 & 0 & 0 \\
Change & $-x$ & $+x$ & $+x$ \\
Equilibrium & $0.100 - x$ & $x$ & $x$ \\
\hline
\end{tabular}
\end{center}

\[ K = \frac{x^2}{0.100 - x} = 1.0\times10^{-5} \]

Because $K$ is very small, very little A reacts, so
$0.100 - x \approx 0.100$:
\[ x^2 \approx (1.0\times10^{-5})(0.100) = 1.0\times10^{-6}
   \qquad x = \mathbf{1.0\times10^{-3}} \]

\textbf{Test it:} $x/[\ce{A}]_0 = 1.0\%$, comfortably under 5\%, so the
approximation is justified. \checkmark
\end{workedexample}

\begin{apctrap}
\textbf{The 5\% check is part of the answer, not a formality.} An FRQ that
asks you to justify an approximation wants the \emph{number} and the
comparison, not the assertion. ``$x$ is small compared with 0.100'' with no
percentage earns nothing.

Rule of thumb: the approximation is safe when $K$ is small relative to the
initial concentration --- roughly $[\ce{A}]_0/K > 400$.
\end{apctrap}

\segment{APPLICATION}{20}{Running the table}

\begin{enumerate}[leftmargin=*,itemsep=0.4em]
  \item $\ce{A <=> 2B}$, $K = 4.0\times10^{-6}$,
        $[\ce{A}]_0 = \SI{0.50}{\Molar}$. Find $[\ce{B}]$ at equilibrium.
        \workspace[2.6cm]
        \keyonly{\begin{apcanswer}$K = (2x)^2/(0.50-x) \approx
        4x^2/0.50$; $4x^2 = 2.0\times10^{-6}$;
        $x = 7.1\times10^{-4}$;
        $[\ce{B}] = 2x = \mathbf{1.4\times10^{-3}}$~M. Check:
        $x/0.50 = 0.14\%$ \checkmark\end{apcanswer}}
  \item Why does a small $K$ make the approximation safe, in one sentence?
        \answerlines[2]{A small $K$ means the equilibrium lies far to the
        left, so only a tiny fraction of the reactant is consumed and
        subtracting $x$ barely changes the initial value.}
\end{enumerate}

\begin{exitticket}
State the two steps that must appear in a complete small-$x$ solution.
\keyonly{\begin{apcanswer}Solve using the approximation, then compute
$x/[\ce{A}]_0$ as a percentage and compare it with 5\%. Omitting the second
step leaves the answer unjustified.\end{apcanswer}}
\end{exitticket}

% =====================================================================
\blocklesson{When the Approximation Fails}{CED 7.7 \textbullet\ Zumdahl \S13.6}

\begin{objectives}
  \item Recognize when the small-$x$ approximation is invalid.
  \item Solve the resulting quadratic or perfect square.
\end{objectives}

\reading{Zumdahl \S13.6}{671--676}

\begin{warmup}[3]
  \item The approximation threshold: \answerblank[1.6cm]{5\%}
  \item ICE stands for: \answerblank[5.8cm]{Initial, Change, Equilibrium}
  \item Large $K$ means the approximation is likely to be:
        \answerblank[2cm]{invalid}
\end{warmup}

\segment{INSTRUCTION A}{25}{The check that saves you}

\notesection{A failed approximation, quantified}{7.7}
\practice{5}

\begin{workedexample}[Worked example 6: the same setup, a bigger $K$]
$\ce{A <=> B + C}$ with $[\ce{A}]_0 = 0.100$ but now $K = 0.20$.

\textbf{Try the approximation:}
$x^2 \approx (0.20)(0.100)$ gives $x = 0.141$.

\textbf{Test it:} $x/[\ce{A}]_0 = 141\%$ --- the ``small'' $x$ is larger
than the entire initial concentration. \textbf{Invalid.}

\textbf{Solve properly:} $x^2 = 0.20(0.100 - x)$, so
$x^2 + 0.20x - 0.020 = 0$ and
\[ x = \frac{-0.20 + \sqrt{0.040 + 0.080}}{2} = \mathbf{0.073} \]

The approximation was wrong by about \textbf{93\%}. This is why the check
is not optional.
\end{workedexample}

\segment{INSTRUCTION B}{20}{Perfect squares}

\notesection{When both sides are squares, take the root}{7.7}
\practice{5}

Some ICE expressions avoid the quadratic formula entirely.

\begin{workedexample}[Worked example 7: hydrogen iodide]
$\ce{H2 + I2 <=> 2HI}$, $K = 50.0$, both reactants starting at
\SI{1.00}{\Molar}.

\begin{center}
\begin{tabular}{@{}lccc@{}}
\hline
 & \ce{H2} & \ce{I2} & \ce{HI} \\
\hline
Initial & 1.00 & 1.00 & 0 \\
Change & $-x$ & $-x$ & $+2x$ \\
Equilibrium & $1.00-x$ & $1.00-x$ & $2x$ \\
\hline
\end{tabular}
\end{center}

\[ K = \frac{(2x)^2}{(1.00-x)(1.00-x)}
     = \left(\frac{2x}{1.00-x}\right)^2 = 50.0 \]

Both sides are perfect squares, so take the square root:
\[ \frac{2x}{1.00-x} = \sqrt{50.0} = 7.07 \]
\[ 2x = 7.07 - 7.07x \quad\Longrightarrow\quad 9.07x = 7.07
   \quad\Longrightarrow\quad x = \mathbf{0.780} \]

So $[\ce{H2}] = [\ce{I2}] = \mathbf{0.220}$ and
$[\ce{HI}] = \mathbf{1.56}$~M.

\textbf{Check:} $(1.56)^2/(0.220)^2 = 50.0$ \checkmark\ --- always worth
30 seconds.
\end{workedexample}

\segment{APPLICATION}{20}{Choosing a method}

\begin{enumerate}[leftmargin=*,itemsep=0.4em]
  \item For each, say whether the small-$x$ approximation is likely valid:
        (i) $K = 10^{-8}$, $[\ce{A}]_0 = 1.0$;
        (ii) $K = 2.0$, $[\ce{A}]_0 = 0.10$;
        (iii) $K = 10^{-6}$, $[\ce{A}]_0 = 10^{-5}$.
        \answerlines[3]{(i) Valid --- $K$ is minute against the initial
        concentration. (ii) Invalid --- $K$ is comparable to $[\ce{A}]_0$.
        (iii) Invalid --- although $K$ is small, so is $[\ce{A}]_0$, and it
        is the \emph{ratio} that matters.}
  \item Solve $\ce{A <=> 2B}$ with $K = 0.50$ and
        $[\ce{A}]_0 = 1.00$~M, using the quadratic.
        \workspace[3cm]
        \keyonly{\begin{apcanswer}$K = 4x^2/(1.00-x) = 0.50$, so
        $4x^2 + 0.50x - 0.50 = 0$.
        $x = [-0.50 + \sqrt{0.25 + 8.0}]/8 = (-0.50+2.872)/8 =
        \mathbf{0.297}$;
        $[\ce{B}] = 0.593$, $[\ce{A}] = 0.703$~M\end{apcanswer}}
\end{enumerate}

\begin{exitticket}
Why is it the \emph{ratio} $[\ce{A}]_0/K$, rather than $K$ alone, that
decides whether the approximation works?
\keyonly{\begin{apcanswer}The approximation requires $x$ to be small
compared with $[\ce{A}]_0$, and $x$ scales with $K$ while the thing it is
compared against is $[\ce{A}]_0$. A small $K$ with an even smaller starting
concentration still gives a large percentage, so only the ratio
settles it.\end{apcanswer}}
\end{exitticket}

% =====================================================================
\blocklesson{Representations and Le Ch\^atelier}{CED 7.8--7.9 \textbullet\ Zumdahl \S13.1, \S13.7}

\begin{objectives}
  \item Interpret particulate and graphical representations of equilibrium.
  \item Predict the effect of a disturbance.
\end{objectives}

\reading{Zumdahl \S13.1, \S13.7}{651--654, 676--683}

\begin{warmup}[3]
  \item Failed approximation $\Rightarrow$ use the:
        \answerblank[2.6cm]{quadratic}
  \item At equilibrium, $Q$ equals: \answerblank[1.6cm]{$K$}
  \item Only what changes $K$? \answerblank[2.6cm]{temperature}
\end{warmup}

\segment{INSTRUCTION A}{25}{Le Ch\^atelier's principle}

\notesection{A system opposes the change imposed on it}{7.9}
\practice{6}

\begin{center}
\renewcommand{\arraystretch}{1.25}
\begin{tabular}{@{}lll@{}}
\hline
\textbf{Disturbance} & \textbf{Shift} & \textbf{Does $K$ change?} \\
\hline
add reactant & right & no \\
add product & left & no \\
remove product & right & no \\
compress (smaller $V$) & toward fewer moles of gas & no \\
expand (larger $V$) & toward more moles of gas & no \\
add inert gas at constant $V$ & \textbf{no shift} & no \\
add a catalyst & \textbf{no shift} & no \\
raise $T$, endothermic forward & right & \textbf{yes, increases} \\
raise $T$, exothermic forward & left & \textbf{yes, decreases} \\
\hline
\end{tabular}
\end{center}

\begin{apctrap}
\textbf{Only temperature changes the value of $K$.} Everything else moves
the system \emph{along} the same curve until $Q$ returns to the unchanged
$K$. A response saying ``adding more reactant increases $K$'' is marked
wrong every time.

Two more that recur: an \textbf{inert gas} added at constant volume changes
the total pressure but not any partial pressure or concentration, so
nothing shifts. And a \textbf{catalyst} speeds both directions equally ---
equilibrium arrives sooner, in the same place.
\end{apctrap}

\notesection{Temperature as a reactant or product}{7.9}

Treat energy as a term in the equation. For an exothermic reaction, write
heat as a \answerblank[2cm]{product}; adding heat then pushes the
equilibrium \answerblank[1.8cm]{left}, exactly as adding any other product
would.

\segment{INSTRUCTION B}{20}{Particulate representations}

\notesection{Reading a picture of an equilibrium}{7.8}
\practice{1}

When a question shows before-and-after particle diagrams:

\begin{itemize}[leftmargin=*,itemsep=0.3em]
  \item Count particles and check that atoms are
        \answerblank[2.4cm]{conserved}.
  \item Compare the ratio of products to reactants against
        \answerblank[1.6cm]{$K$}.
  \item Two diagrams showing \emph{no change in counts} between successive
        times indicate the system has reached
        \answerblank[2.6cm]{equilibrium}.
\end{itemize}

\segment{APPLICATION}{20}{Predicting shifts}

For $\ce{N2(g) + 3H2(g) <=> 2NH3(g)}$, $\Delta H = -92$~kJ:

\begin{enumerate}[leftmargin=*,itemsep=0.35em]
  \item Add \ce{N2}: \answerblank[3cm]{shifts right}
  \item Remove \ce{NH3}: \answerblank[3cm]{shifts right}
  \item Compress the vessel: \answerblank[5.4cm]{right --- 4 mol gas to 2}
  \item Add argon at constant volume:
        \answerblank[3cm]{no shift}
  \item Raise the temperature:
        \answerlines[2]{Shifts left, and $K$ decreases. The forward reaction
        is exothermic, so heat behaves as a product; adding heat pushes the
        system back toward reactants.}
  \item Add a catalyst: \answerblank[5.8cm]{no shift; equilibrium sooner}
\end{enumerate}

\begin{exitticket}
Name the only disturbance that changes the value of $K$, and explain why
the others do not.
\keyonly{\begin{apcanswer}Temperature. The others move concentrations or
pressures, which changes $Q$; the system then shifts until $Q$ returns to
the same $K$. Temperature instead changes the constant itself, so the
system settles at a genuinely different ratio.\end{apcanswer}}
\end{exitticket}

% =====================================================================
\blocklesson{$Q$ versus $K$ Quantitatively}{CED 7.10 \textbullet\ Zumdahl \S13.7}

\begin{objectives}
  \item Calculate $Q$ and compare it with $K$ to predict direction.
  \item Justify a Le Ch\^atelier prediction quantitatively.
\end{objectives}

\reading{Zumdahl \S13.7}{676--683}

\begin{warmup}[3]
  \item Adding a product shifts the equilibrium:
        \answerblank[1.8cm]{left}
  \item An inert gas at constant volume causes:
        \answerblank[2.4cm]{no shift}
  \item Compressing shifts toward:
        \answerblank[4cm]{fewer moles of gas}
\end{warmup}

\segment{INSTRUCTION A}{25}{The calculation behind the principle}

\notesection{$Q$ is the answer to ``which way?''}{7.10}
\practice{5}

Le Ch\^atelier tells you the direction; $Q$ \emph{proves} it. On an FRQ,
the quantitative justification is usually worth the point that the
qualitative one is not.

\begin{workedexample}[Worked example 8: predicting with $Q$]
For $\ce{H2 + I2 <=> 2HI}$, $K = 4.5$. A vessel is charged with
$[\ce{H2}] = 0.40$, $[\ce{I2}] = 0.30$, $[\ce{HI}] = 0.40$~M.

\[ Q = \frac{(0.40)^2}{(0.40)(0.30)} = \frac{0.16}{0.12} = 1.3 \]

$Q = 1.3 < K = 4.5$, so there are \emph{too few products}: the reaction
proceeds \textbf{to the right} until $Q$ rises to 4.5.

Note that this required no Le Ch\^atelier reasoning at all --- the
comparison settles it directly, and it works even for a fresh mixture that
was never at equilibrium.
\end{workedexample}

\begin{apctrap}
Le Ch\^atelier applies to \emph{disturbing a system already at
equilibrium}. $Q$ versus $K$ works for \textbf{any} mixture, equilibrium or
not. When a question gives you arbitrary starting concentrations, $Q$ is the
tool --- Le Ch\^atelier has nothing to say about a mixture that was never at
equilibrium in the first place.
\end{apctrap}

\segment{APPLICATION}{20}{Direction from data}

\begin{enumerate}[leftmargin=*,itemsep=0.4em]
  \item $K = 4.5$ for $\ce{H2 + I2 <=> 2HI}$. For each mixture, calculate
        $Q$ and state the direction:
        \begin{enumerate}[leftmargin=1.4em,itemsep=0.2em]
          \item $[\ce{H2}]=[\ce{I2}]=[\ce{HI}]=0.50$
                \answerblank[5cm]{$Q = 1.0$; shifts right}
          \item $[\ce{H2}]=0.10$, $[\ce{I2}]=0.10$, $[\ce{HI}]=0.30$
                \answerblank[5cm]{$Q = 9.0$; shifts left}
        \end{enumerate}
  \item A system at equilibrium has extra product added. Explain, using
        $Q$, why it shifts left. \answerlines[3]{Adding product increases
        the numerator of $Q$ without changing $K$, so $Q$ becomes greater
        than $K$. A system with $Q > K$ has too many products relative to
        equilibrium, so the reverse reaction dominates until $Q$ falls back
        to $K$.}
\end{enumerate}

\begin{exitticket}
When is $Q$ the right tool and Le Ch\^atelier's principle the wrong one?
\keyonly{\begin{apcanswer}When the mixture was never at equilibrium --- for
instance an arbitrary set of starting concentrations. Le Ch\^atelier
describes the response to disturbing an existing equilibrium; comparing $Q$
with $K$ works for any mixture at all.\end{apcanswer}}
\end{exitticket}

% =====================================================================
\blocklesson{Solubility Equilibria}{CED 7.11--7.12 \textbullet\ Zumdahl \S16.1}

\begin{objectives}
  \item Relate $K_{sp}$ to molar solubility.
  \item Explain and calculate the common-ion effect.
\end{objectives}

\reading{Zumdahl \S16.1}{805--812}

\begin{warmup}[3]
  \item $Q < K$ means the reaction shifts:
        \answerblank[1.8cm]{right}
  \item Pure solids appear in $K$? \answerblank[1.6cm]{no}
  \item Adding a product shifts equilibrium:
        \answerblank[1.8cm]{left}
\end{warmup}

\segment{INSTRUCTION A}{25}{The same machinery, applied to a solid}

\notesection{$K_{sp}$ is just $K$}{7.11}
\practice{5}

A saturated solution is a system at equilibrium with undissolved solid:
\[ \ce{CaF2(s) <=> Ca^2+(aq) + 2F^-(aq)} \qquad
   K_{sp} = [\ce{Ca^2+}][\ce{F^-}]^2 \]

The solid is omitted for the reason you already know --- it is a
\answerblank[2.4cm]{pure solid}.

\begin{center}
\begin{tabular}{@{}lll@{}}
\hline
\textbf{Type} & \textbf{$K_{sp}$ in terms of $s$} & \textbf{Solubility} \\
\hline
\ce{MX} & $s^2$ & $s = \sqrt{K_{sp}}$ \\
\ce{MX2} or \ce{M2X} & $4s^3$ & $s = \sqrt[3]{K_{sp}/4}$ \\
\hline
\end{tabular}
\end{center}

\begin{apctrap}
\textbf{$K_{sp}$ is a constant; solubility is a position.} $K_{sp}$ has one
value at a given temperature; the solubility changes with whatever else is
in the water.

And \textbf{$K_{sp}$ values may be compared directly only for salts
producing the same number of ions.} \ce{CaF2}
($K_{sp} = 4.0\times10^{-11}$) is far more soluble than \ce{BaCrO4}
($8.5\times10^{-11}$) despite the smaller constant, because $s$ enters to
different powers.
\end{apctrap}

\segment{INSTRUCTION B}{20}{The common-ion effect}

\notesection{Le Ch\^atelier, one last time}{7.12}
\practice{6}

Dissolve a salt in a solution already containing one of its ions and you
have added a \answerblank[2cm]{product}. The equilibrium shifts
\answerblank[1.8cm]{left}, so \emph{less} dissolves --- while $K_{sp}$
stays exactly the same.

\begin{workedexample}[Worked example 9: silver chloride]
$K_{sp}(\ce{AgCl}) = 1.6\times10^{-10}$.

\textbf{In pure water:} $s = \sqrt{1.6\times10^{-10}} =
\mathbf{1.3\times10^{-5}}$~M.

\textbf{In \SI{0.10}{\Molar} \ce{NaCl}:} the chloride is essentially all
from the \ce{NaCl}, so
\[ s = \frac{K_{sp}}{[\ce{Cl^-}]} = \frac{1.6\times10^{-10}}{0.10}
   = \mathbf{1.6\times10^{-9}}~\text{M} \]

Suppressed by a factor of about \textbf{7900} --- and $K_{sp}$ never moved.
\end{workedexample}

\segment{APPLICATION}{20}{Solubility calculations}

\begin{enumerate}[leftmargin=*,itemsep=0.4em]
  \item Calculate the molar solubility of \ce{CaF2}
        ($K_{sp} = 4.0\times10^{-11}$) in pure water.
        \workspace[2.2cm]
        \keyonly{\begin{apcanswer}$4s^3 = 4.0\times10^{-11}$;
        $s^3 = 1.0\times10^{-11}$;
        $s = \mathbf{2.2\times10^{-4}}$~M\end{apcanswer}}
  \item Calculate it in \SI{0.025}{\Molar} \ce{NaF}.
        \workspace[2.2cm]
        \keyonly{\begin{apcanswer}$[\ce{F-}] \approx 0.025$;
        $4.0\times10^{-11} = s(0.025)^2$;
        $s = \mathbf{6.4\times10^{-8}}$~M\end{apcanswer}}
  \item State what happened to $K_{sp}$ between (a) and (b), and why.
        \answerlines[2]{Nothing --- $K_{sp}$ depends only on temperature.
        Only the equilibrium position moved, with much more fluoride and
        much less calcium in solution.}
\end{enumerate}

\begin{exitticket}
Explain the common-ion effect using $Q$ rather than Le Ch\^atelier.
\keyonly{\begin{apcanswer}Adding the common ion raises the ion product $Q$
above $K_{sp}$. A system with $Q > K$ shifts toward the reactant side ---
here, toward the undissolved solid --- until $Q$ falls back to $K_{sp}$,
which leaves less of the salt dissolved.\end{apcanswer}}
\end{exitticket}

\end{document}
